%% Chapter 8 — Section 8.4: Exercises
%% Source: PDF pages 109–110 (book pages 96–97)

\section{Exercises}
\label{sec:8.4}

\begin{enumerate}

\item[8.1] Prove Corollary~8.8.

\item[8.2] Implement HMC to sample a bivariate Gaussian target $\mathcal{N}(\mu, \Sigma)$.
Consider parameter values $\mu = (0,0)$, $\sigma_X = 1$, $\sigma_Y = 1$ and two different
values of the correlation $\rho = 0$ and $\rho = 0.8$. Provide plots showing:
\begin{enumerate}
  \item The value of the Hamiltonian along trajectories of the leapfrog integrator for various
    choices of the step-size $\epsilon$. The goal is to show that if $\epsilon$ is small these
    trajectories are close to level-sets from the Gaussian target, but the accuracy deteriorates
    if $\epsilon$ is too large.
  \item A scatterplot of samples from HMC with different choices of step-size $\epsilon$ and
    duration parameter $\lambda$.
\end{enumerate}

\item[8.3] Recall from Exercise~6.3 that a response random variable $Y$ given explanatory
variables $(X_1, \ldots, X_d)$ follows a logistics regression model if
\[
Y = \begin{cases} 1 & \text{with probability}\quad \alpha, \\ 0 & \text{with probability}\quad 1 - \alpha, \end{cases}
\]
where
\[
\alpha = \frac{1}{1 + \exp\!\left(-(w_0 + w_1 X_1 + \ldots + w_d X_d)\right)}.
\]
\begin{enumerate}
  \item \textbf{Setting and Preparation}: Same as in Exercise~6.3.
  \item \textbf{Hamiltonian Monte Carlo}: Adopt a Bayesian formulation, specifying a Gaussian
    prior $w \sim \mathcal{N}(0, \sigma^2 I)$.
  \begin{enumerate}
    \item[(b)] Based on the derivations in Exercise~6.3, write down the potential $V(w)$ needed
      to implement HMC.
    \item[(c)] Choose
      \[
        M = \begin{bmatrix} 10 & 0 & 0 \\ 0 & 1 & 0 \\ 0 & 0 & 1 \end{bmatrix}
      \]
      and kinetic energy $K(p) = p^T M p$. Write down the leapfrog integrator applied to the
      Hamiltonian $H(w,p) = V(w) + K(p)$. Plot leapfrog trajectories for different choices of
      step-size.
    \item[(d)] Implement HMC for posterior sampling, tuning the step-size and duration
      parameters. Plot the value of the computed posterior mean of $w$ against the number of
      HMC samples.
  \end{enumerate}
\end{enumerate}

\item[8.4] Consider a 100-dimensional multivariate Gaussian with independent coordinates, mean
vector $\mu = 0$, and variance vector $\sigma^2 = (0.01, 0.02, \ldots, 1)$.
\begin{enumerate}
  \item[(a)] Sample using RWMH and plot histograms for the coordinates with variance $0.98$,
    $0.99$, and $1$.
  \item[(b)] Sample using MALA tuning the step-size. Plot histograms for the coordinates with
    variance $0.98$, $0.99$, and $1$.
  \item[(c)] Sample using HMC with $\lambda = 150$ and tuning the step-size. Plot histograms
    for the coordinates with variance $0.98$, $0.99$, and $1$.
  \item[(d)] For each of the 100 coordinates, compute the sample mean and sample variance of
    the samples you obtained with RWMH, MALA, and HMC. For each method, draw scatter plots of
    the 100 sample means and variances.
\end{enumerate}

\end{enumerate}
